%!TEX root = ../hutbthesis_main.tex
% 设置英文摘要
% 说明：本摘要为开题阶段占位框架，须与中文摘要保持一致，实验数据待系统测试通过后据实补写。
\keywordsen{Underwater Simulation, Sensor Modeling, Underwater Acoustic Communication Link, Sonar Imaging, Delay and Noise, Operation Training}
\begin{abstracten}

In scenarios such as marine engineering, underwater detection, underwater search and rescue, and underwater operation training, the underwater environment is characterized by high noise, high signal delay, obvious carrier drift, severe occlusion, and low visibility, which significantly degrade the perception quality of optical and acoustic sensors. Meanwhile, the underwater acoustic communication link suffers from narrow bandwidth, strong multipath effect, long propagation delay, and packet loss. Collecting data in real underwater environments is costly, risky, and difficult to reproduce. This thesis focuses on the simulation of sensors and communication links in underwater working environments. Based on a simulation platform, an underwater sensor and communication link simulation system is built to model and visualize environmental factors such as underwater noise, delay, drift, occlusion, low visibility, and sonar imaging, and to simulate the sensor data and the transmission process of the communication link, so as to provide a repeatable and low-cost experimental environment for perception algorithms, communication strategies, and underwater operation training.

The main research contents include: modeling of the underwater working environment and parameter configuration; perception simulation of optical and acoustic sensors under noise, delay, drift, occlusion, and low-visibility conditions; simulation of communication link characteristics such as delay, bandwidth, packet loss, and multipath; simulation and display of sonar imaging results; and interactive testing and evaluation oriented to algorithm verification and operation training. All functions will undergo functional and performance tests in the specified simulation scenes, and the experimental results and quantitative metrics will be supplemented after the system is implemented and verified.

\end{abstracten}
